We estimate what would happen if a group of units (customers, patients, calls) were served by a channel other than the one serving them today, without a randomized experiment. We cast identification as a change of measure: the counterfactual mean is an integral over the donor population weighted by the Radon--Nikodym derivative, which a calibrated classifier identifies up to a constant. The focus is a problem that is common in operational data and rarely formalized: each channel records covariates in its own way. A variable recorded in only one channel breaks support by construction, and a confident classifier ($p=1$) reveals it; dropping it is safe when it does not predict the outcome given the other covariates. A variable measured with different instruments in each channel is more dangerous: it can preserve support, and the estimator then converges to a shifted value that no analysis of the observed data can distinguish from the truth. We also show that, without support, the doubly robust estimator implicitly extrapolates the outcome model, and we adapt sensitivity bounds for unobserved confounding. When the target group is a cutoff on a known score, the weights follow from the ordinary channel propensity with the cutoff imposed; estimating it with all units of the premium channel reduces the dispersion by up to 59\,\% in our simulation when overlap is good. A simulation with predictions registered before running it tests the main results and also finds, in our design, that the unit-level bootstrap overstates dispersion when the estimated weights are noisy. With public data (self-reported versus measured body mass index in NHANES), the predicted bias of differential measurement is approximately confirmed, while the overlap diagnostics improve rather than flag it.
